Una massa esfèrica d'acer de 0,300~kg està subjecta a una vareta metàl·lica prima i de massa negligible. Aquesta vareta està clavada verticalment a una massa fixa, de manera que l'extrem on hi ha la massa pot oscil·lar lliurement. Si apliquem una força de 8,00~N sobre l'esfera, aquesta es desplaça 4,0~cm. Suposeu que aquest desplaçament és rectilini i horitzontal, com mostra la figura, i que la força recuperadora de la vareta obeeix la llei de Hooke.

\figura[0.4]{fig1.png}

\begin{apartat}{1}
Calculeu la constant elàstica $k$. Deduïu, a partir de la segona llei de Newton, la fórmula per a obtenir la freqüència d'oscil·lació i calculeu el període d'oscil·lació.
\end{apartat}

\begin{apartat}{1}
Deduïu, a partir de l'equació del moviment harmònic simple (MHS), la fórmula per a obtenir l'acceleració màxima i calculeu-la en aquest cas.
\end{apartat}
